\chapter{Ersatzschaltbilder der Netztechnik}
\label{chap:06-ersatzschaltbilder}

%\chapter{Zusammengesetzte Ersatzschaltbilder und Netzreduktion}
%\label{chap:ersatznetze}

\section{Ziel der Netzmodellierung}

Ein elektrisches Netz besteht aus Quellen, Transformatoren, Leitungen,
Kabeln, Lasten und gegebenenfalls rotierenden Maschinen oder Umrichtern.
Für eine konkrete Berechnung wird daraus ein Modell gebildet, das die für
die jeweilige Fragestellung relevanten Eigenschaften erhält.

Das Modell eines Kurzschlusses muss nicht dieselben Details enthalten wie
ein Modell für Spannungsfall, Verlustleistung oder Erdschlussstrom. Die
Modellierung beginnt deshalb mit der Frage, welche Größe bestimmt werden soll.

\section{Impedanzen auf eine gemeinsame Spannungsebene beziehen}

Für einen idealen Transformator mit Übersetzungsverhältnis
\[
  a=\frac{U_1}{U_2}
\]
wird eine sekundärseitige Impedanz auf die Primärseite bezogen durch
\[
  Z_2'=a^2Z_2.
\]
Umgekehrt wird eine primärseitige Impedanz auf die Sekundärseite bezogen durch
\[
  Z_1'=\frac{Z_1}{a^2}.
\]

Diese Regeln gelten für die Impedanzen der jeweils gekoppelten Ersatzkreise.
Bei Drehstromtransformatoren sind die Spannungs- und Strangdefinitionen
konsistent zu halten.

\subsection{Beispiel}

Eine Impedanz auf der 0{,}4-kV-Seite eines idealisierten Transformators
soll auf die 20-kV-Seite bezogen werden. Mit dem Verhältnis
\[
  a=\frac{20}{0{,}4}=50
\]
gilt
\[
  Z'=50^2 Z=2500Z.
\]

Eine sekundärseitige Impedanz von \(0{,}01\,\Omega\) entspricht damit
auf der 20-kV-Seite einer Impedanz von
\[
  Z'=25\,\Omega.
\]

Die Umrechnung verändert nicht das physikalische Verhalten des Netzes.
Sie verändert nur die Darstellung der Impedanz auf einer anderen
Spannungsseite.

\section{Komplexe Impedanzen addieren}

Für eine Reihenschaltung von Impedanzen gilt
\[
  Z_{\mathrm{ges}}=\sum_{i=1}^{n}Z_i.
\]

Mit
\[
  Z_i=R_i+jX_i
\]
folgt
\[
  R_{\mathrm{ges}}=\sum_i R_i,\qquad
  X_{\mathrm{ges}}=\sum_i X_i.
\]
Erst anschließend wird der Betrag bestimmt:
\[
  |Z_{\mathrm{ges}}|=
  \sqrt{R_{\mathrm{ges}}^2+X_{\mathrm{ges}}^2}.
\]

Im Allgemeinen gilt nicht
\[
  |Z_1+Z_2|=|Z_1|+|Z_2|.
\]
Die Beträge einzelner komplexer Impedanzen dürfen deshalb nicht einfach
addiert werden.

Für eine Parallelschaltung gilt
\[
  \frac{1}{Z_{\mathrm{ges}}}=\sum_i\frac{1}{Z_i}.
\]
Es ist häufig zweckmäßig, zuerst die Admittanzen
\[
  Y_i=\frac{1}{Z_i}
\]
zu bilden und diese zu addieren.

\section{Thévenin-Ersatzquelle}

Ein lineares Netzwerk kann aus Sicht eines betrachteten Anschlusses unter
geeigneten Voraussetzungen durch eine Ersatzspannungsquelle und eine
Serienimpedanz dargestellt werden.

Die Thévenin-Ersatzspannung \(U_{\mathrm{Th}}\) ist die Leerlaufspannung
am betrachteten Anschluss. Die Thévenin-Impedanz \(Z_{\mathrm{Th}}\)
ist die von diesem Anschluss aus gesehene Impedanz bei deaktivierten
unabhängigen Quellen. Ideale Spannungsquellen werden dabei kurzgeschlossen,
ideale Stromquellen geöffnet. Abhängige Quellen müssen bei der Bestimmung
der Ersatzimpedanz weiterhin berücksichtigt werden.

Bei der Kurzschlussberechnung wird die Ersatzimpedanz aus Sicht der
Fehlerstelle gebildet. Dadurch lässt sich ein komplexes Netz auf ein
einfaches Ersatzmodell reduzieren.

\section{Kurzschlussstrom im Ersatznetz}

Für einen einphasigen Ersatzkreis mit Fehlerimpedanz \(Z_f\) gilt
\[
  I_k=\frac{U_{\mathrm{Th}}}{Z_{\mathrm{Th}}+Z_f}.
\]

Bei einem symmetrischen dreipoligen Fehler wird das Mitsystem-Ersatznetz
verwendet. Die Spannung ist dann die Strangspannung:
\[
  U_{\mathrm{ph}}=\frac{U_{LL}}{\sqrt{3}}.
\]
Daher gilt für den Betrag des symmetrischen Kurzschlussstroms
\[
  I_k=\frac{U_{LL}}{\sqrt{3}|Z_{\mathrm{Th}}+Z_f|}.
\]

Diese Gleichung beschreibt das vereinfachte Ersatznetz. Bei normgerechten
Kurzschlussberechnungen können je nach Berechnungsziel zusätzliche Faktoren,
Anfangsbedingungen, Motorbeiträge, Spannungsfaktoren und Korrekturen
notwendig sein.

\section{Per-Unit-System als einheitliche Darstellung}

Mit einer gewählten dreiphasigen Bezugsleistung \(S_B\) und einer
Leiter-Leiter-Bezugsspannung \(U_B\) werden die Bezugsgrößen definiert als
\[
  I_B=\frac{S_B}{\sqrt{3}U_B},
  \qquad
  Z_B=\frac{U_B^2}{S_B}.
\]

Der Per-Unit-Wert einer Impedanz ist
\[
  z_{\mathrm{pu}}=\frac{Z}{Z_B}.
\]

Für einen Wechsel der Bezugsbasis gilt
\[
  z_{\mathrm{pu,new}}
  =z_{\mathrm{pu,old}}
   \frac{S_{B,\mathrm{new}}}{S_{B,\mathrm{old}}}
   \left(\frac{U_{B,\mathrm{old}}}{U_{B,\mathrm{new}}}\right)^2.
\]

Werden die Bezugsspannungen über ideale Transformatoren konsistent gewählt,
bleiben Per-Unit-Impedanzen auf beiden Seiten häufig besonders einfach
vergleichbar. Die Bezugsspannungen müssen dabei dem Übersetzungsverhältnis
entsprechen.

\section{Durchgängiges Rechenbeispiel}

Gegeben sei ein vereinfachtes Netz mit:
\begin{itemize}
  \item Einspeisung mit Kurzschlussleistung \(S_k''=250\,\mathrm{MVA}\)
        auf der 20-kV-Seite;
  \item Transformator \(20/0{,}4\,\mathrm{kV}\);
  \item Kabel auf der 0{,}4-kV-Seite;
  \item dreipoligem Kurzschluss am Kabelende.
\end{itemize}

Die Netzimpedanz wird zunächst auf der 20-kV-Seite abgeschätzt:
\[
  |Z_{\mathrm{Netz}}|
  \approx\frac{U^2}{S_k''}
  =\frac{(20\,\mathrm{kV})^2}{250\,\mathrm{MVA}}
  =1{,}6\,\Omega.
\]

Diese Formel liefert eine Näherung für den Betrag der Netzimpedanz bei
der betrachteten Spannungsebene. Sie bestimmt nicht deren Wirk- und
Blindanteil getrennt.

Für die Übertragung auf die 0{,}4-kV-Seite ist das Spannungsverhältnis
\[
  a=\frac{20}{0{,}4}=50.
\]
Damit ergibt sich näherungsweise
\[
  |Z_{\mathrm{Netz}}'|=\frac{1{,}6}{50^2}
  =0{,}00064\,\Omega.
\]

Anschließend werden die auf die Niederspannungsseite bezogene Netzimpedanz,
die Transformatorimpedanz und die Kabelimpedanz komplex addiert:
\[
  Z_{\mathrm{ges}}
  =Z_{\mathrm{Netz}}'+Z_{\mathrm{Tr}}+Z_{\mathrm{Kabel}}.
\]

Der anfängliche symmetrische Kurzschlussstrom ergibt sich im vereinfachten
Ersatzmodell zu
\[
  I_k''=\frac{U_{LL}}{\sqrt{3}|Z_{\mathrm{ges}}|}.
\]

Ein numerischer Endwert ist erst dann fachlich belastbar, wenn die
komplexen Transformator- und Kabelimpedanzen sowie die verwendeten
Kurzschlussannahmen vollständig bekannt sind.

\section{Plausibilitätsprüfungen}

Nach einer Netzberechnung sollten mindestens folgende Prüfungen erfolgen:
\begin{itemize}
  \item Stimmen Spannungs-, Strom- und Impedanzdimensionen?
  \item Wurden alle Impedanzen auf dieselbe Spannungsebene bezogen?
  \item Wurden komplexe Impedanzen und nicht nur ihre Beträge addiert?
  \item Ist der errechnete Kurzschlussstrom kleiner als die durch die
        Ersatzimpedanz begrenzte plausible Größenordnung?
  \item Sind die Beiträge von Motoren, Generatoren und Umrichtern
        im verwendeten Modell berücksichtigt oder bewusst vernachlässigt?
  \item Entspricht die verwendete Spannung der Definition als Strang-
        oder Leiter-Leiter-Spannung?
\end{itemize}

\section{Übungsaufgaben}

\begin{enumerate}
  \item Beziehen Sie eine Impedanz von \(0{,}05\,\Omega\) auf der
        0{,}4-kV-Seite auf die 20-kV-Seite.
  \item Addieren Sie \(Z_1=0{,}1+j0{,}3\,\Omega\) und
        \(Z_2=0{,}2+j0{,}1\,\Omega\). Bestimmen Sie Betrag und Winkel
        der Gesamtimpedanz.
  \item Erklären Sie die Bedeutung der Thévenin-Ersatzimpedanz für die
        Berechnung eines Kurzschlussstroms.
  \item Welche zusätzlichen Informationen sind erforderlich, um das
        durchgerechnete Netzbeispiel mit einem belastbaren Zahlenwert
        abzuschließen?
\end{enumerate}